\subsection{正交函数组}
\idx{正交函数组}{zhengjiaohanshuzu}
\begin{definition}
  \label{def:orthogonal_function}
  设$f(x), g(x) $为定义在区间$[a,b]$上的函数, 积分
  \begin{equation}
    (f, g) \equiv \int_a^b f(x) g(x) dx
  \end{equation}
叫做二函数$f(x)$和$g(x)$的内积$(f,g)$或$(fg)$.
\end{definition}
内积满足施瓦茨不等式(Schwarz inequality)
\begin{equation}
  |(f,g)|^2 \leq (f,f)(g,g)
\end{equation}
当且仅当$f(x)$和$g(x)$线性相关时,等号成立.
如果$(f,g) = 0$,则称$f(x)$和$g(x)$是\textbf{正交}的. 函数$f(x)$和它自己的内积称为\textbf{范数}或\textbf{模},记作
$N f$,
\begin{equation}
  N f \equiv (f,f) = \int  f(x)^2 dx
\end{equation}
范数为1的函数称为\textbf{单位函数}或\textbf{归一化函数}. 一组归一化函数$\varphi_1(x), \varphi_2(x), \cdots, \varphi_n(x)$,其中任意二者皆正交，
称\textbf{正交归一组}, 表示为
\begin{equation}
  (\varphi_i, \varphi_j) = \delta_{ij} .
\end{equation}
对于复变函数， 内积定义为
\begin{equation}
  ( f, g ) = \int  f^*(x) g(x) dx
\end{equation}
其中$\bar f$为$f$的复共轭.  若$N f = \int | f|^2 dx = 1$, 称函数$f(x)$是归一化的. 

若对于所有$x$都成立的常系数齐次线性关系
\begin{equation}
  \sum_{i=1}^n c_i \varphi_i(x) = 0,
\end{equation}
其中系数不全为零, 我们称这$n$个函数$f_1, f_2, \cdots, f_n$是\textbf{线性相关}的. 否则称为\textbf{线性无关}的. 值得注意的是
正交组中的各函数永远是线性无关的. 因为, 若恒等式
\begin{equation}
  \sum_{i=1}^n c_i \varphi_i(x) = 0,
\end{equation}
成立，我们可以左乘$\varphi_j(x)$并积分可得$c_j = 0$.

我们可以类比于线性代数中向量的施密特正交化的方法对函数组进行正交化. 设$f_1(x), f_2(x), \cdots, f_n(x)$是线性无关的函数组, 我们可以通过
\begin{equation}
  \varphi_1(x) = f_1(x), \quad \varphi_2(x) = f_2(x) - \frac{(f_2, \varphi_1)}{(\varphi_1, \varphi_1)} \varphi_1(x), \cdots
  \varphi_n(x) = f_n(x) - \sum_{j=1}^{n-1} \frac{(f_n, \varphi_j)}{(\varphi_j, \varphi_j)} \varphi_j(x)
\end{equation}

\begin{example}
在 $[0,1]$ 上把 $\{1,x\}$ 正交化.
\end{example}
\begin{solution}
取 $\varphi_1=1$,则 $(\varphi_1,\varphi_1)=1$,$(x,1)=\int_0^1 x\,\mathrm{d}x=1/2$,所以
\begin{equation}
  \varphi_2 = x - \frac12 .
\end{equation}
检验 $\int_0^1 (x-1/2)\,\mathrm{d}x=0$.其范数平方是 $\int_0^1 (x-1/2)^2\,\mathrm{d}x=1/12$,对应的归一化函数是 $\sqrt{12}\,(x-1/2)$.
在对称区间 $[-1,1]$ 上, $(x,1)$ 本来就是零, $1$ 和 $x$ 已经正交,不必再减.
\end{solution}

下面我们来讨论正交函数组的完备性. 设$\varphi_1(x), \varphi_2(x), \cdots, \varphi_n(x)$是正交归一组, 我们可以定义
\begin{equation}
  f(x) = \sum_{i=1}^n c_i \varphi_i(x)
\end{equation}
其中系数$c_i$为
\begin{equation}
  c_i = (f, \varphi_i) = \int f(x) \varphi_i(x) dx
\end{equation}
成为\textbf{展开系数}或\textbf{分量}.

由关系式
\begin{equation}
  \int \left( f(x) - \sum_{i=1}^n c_i \varphi_i(x) \right)^2 dx \geq  0
\end{equation}
展开个平方并逐项积分, 我们可以得到
\begin{equation}
 \int f^2(x) dx +  \sum_{i=1}^n c_i^2  -  2 \sum_{i=1}^n   c_i \int f  \varphi_i  dx \geq 0 
\end{equation}
因此有,
\begin{equation}
  \sum_{i=1}^n c_i^2 \leq N f 
\end{equation}
上式可以令$n \to \infty$, 也可以得到
\begin{equation}
  \sum_{i=1}^{\infty} c_i^2 \leq N f
\end{equation}
这称为\textbf{贝塞尔不等式}(Bessel's inequality). 它对每一个正交归一组都成立, 它证明了展开系数的平方和是收敛有界的.

假如有一正交归一组$\varphi_1(x), \varphi_2(x), \cdots, $, 使得对任意分段连续函数$f(x)$,都有$n$使
平均平方误差
$$
\int \left( f(x) - \sum_{i=1}^n c_i \varphi_i(x) \right)^2 dx 
$$
小于任意给定正数$\epsilon$, 那么我们称$\varphi_1(x), \varphi_2(x), \cdots$是\textbf{完备的}. 对于完备的正交归一组而言, 贝塞尔不等式
对任意函数$f(x)$都成为等式:
\begin{equation}
  \sum_{i=1}^{\infty} c_i^2 = N f
  \label{eq:completeness}
\end{equation}
这个关系称为``完备性关系''. 更一般的形式为
\begin{equation}
  \sum_{i=1}^{\infty} c_i d_i = (f, g) \quad \text{其中} \quad  c_i = (f, \varphi_i) , d_i = (g, \varphi_i)
\end{equation}
可以验证, 对$f+g$ 应用完备性关系~\ref{eq:completeness},
\begin{equation}
  N (f + g) = Nf  + Ng + 2 (f, g) = \sum_{i = 1}^{\infty} (c_i + d_i)^2 = \sum_{i=1}^{\infty} (c_i^2 + d_i^2 + 2  c_i d_i) 
\end{equation}
并减去相应的$f,g$的范数.

%  该不等式说明了展开系数的平方和是有界的, 但并不能保证展开系数的平方和等于函数的模平方. 也就是说, 我们并不能保证

\subsection{傅里叶级数的引入}
\label{subsec:fourier_series}
Fourier曾试图解决一个问题,在此过程中发展了傅里叶级数的概念.此问题处理长度为$\ell$的一维均匀棒的热传播过程,即
给定初始温度分布,求解在$t$时刻,$x$处的温度,$T(x,t)$.该温度场满足微分方程
$$
\frac{\partial^2}{\partial x^2} T(x,t) = \frac{\partial}{\partial t} T(x,t), \quad  T(0,t) = T(\ell,t) = 0 .  
$$
该微分方程的一特殊解可以由分离变量法得到,具体方法后面会学习.令$T(x,t) = f(x) g(t)$,代入该方程有
$$
  f''(x) g(t) = f(x) g'(t) ,
$$
两边同除以$f(x)g(t)$得到
$$
\frac{f''(x)}{f(x)} = \frac{g'(t)}{g(t)}   .
$$
由于左式是与$t$无关的,而右边是与$x$无关的,于是唯一可能相等的情况为二者皆为同一个常数$c$.
这时候若$c>0$,可知$\lim_{t\to \infty} g(t) = \infty$,因此我们令$c=-k^2$.
那么可以知道
$$
f(x) =  A \sin{k x} + B \cos{k x},    
$$
由边界条件确定$B = 0$,且$k = \frac{n\pi}{\ell}$.求解$g(t)$得 $g(t) = e^{-k^2 t}$.
该方程的解可以表示为
$$
T(x,t) = A \sin{\left( \frac{n\pi}{\ell} x \right)} e^{-\frac{n^2\pi^2}{\ell^2} t}  
$$
其中$A$由其他条件确定.
可以知道$t=0$时,$T$为一个正弦函数,节点数由$n$决定.但是,更普遍的情况是起始时刻的温度分布并不是正弦函数.
那该如何决定稍后某时刻$t$的分布呢?傅里叶发现我们并不需要对此问题反反复复求解.
对于线性微分方程可以知道若$T_1(x,t)$和$T_2(x,t)$满足微分方程
$$
    \frac{\partial^2}{\partial x^2} T(x,t) = \frac{\partial}{\partial t} T(x,t)
$$
那么二者的任意线性组合
$$
T_3 = \alpha T_1  + \beta T_2  
$$
也是其解.
傅里叶得到结论,对于任意一种光滑的温度分布,其解总可以写成一系列有着不同振幅和模式的正弦函数的线性叠加,即
$$
  T(x,t) = \sum_{n=1}^{\infty} A_n \sin {\left( \frac{n\pi}{\ell} x \right)} e^{-\frac{n^2\pi^2}{\ell^2} t} .     
$$
若初始温度分布函数为$f(x) = T(x,0)$,那么
$$
  f(x) =      \sum_{n=1}^{\infty} A_n \sin {\left( \frac{n\pi}{\ell} x \right)}
$$
上式什么样的函数可以表示为正弦和余弦函数的叠加呢?每个模式的振幅$A_n$该如何计算呢?
这就利用了正余弦函数的特征.

记
\begin{equation}
 (f, g) = \int_0^{2\pi} f^{*}(x) g(x) dx\,,
\end{equation}

可以验证,积分区间取 $[0,2\pi]$ 时
$$
(\sin{ n x }, \sin{m x })  =  (\cos{ n x } \cos {m x })  =  \delta_{mn}\pi  \; (m\neq 0, n\neq 0)
$$

$$
(\sin{ n x },  \cos{m x })  =  0
$$

\begin{example}
用积化和差核对 $[0,2\pi]$ 上的正交性,再换到热棒的区间 $[0,\ell]$.
\end{example}
\begin{solution}
$\sin A\sin B=\frac12\bigl[\cos(A-B)-\cos(A+B)\bigr]$.$m,n$ 是正整数且 $m\neq n$ 时,两个余弦都在 $[0,2\pi]$ 上积满整数个周期,积分是零.$m=n$ 时
\begin{equation}
  \int_0^{2\pi}\frac12\bigl(1-\cos 2nx\bigr)\,\mathrm{d}x=\pi.
\end{equation}
余弦与余弦同样.正弦与余弦化成 $\sin(n+m)x$ 和 $\sin(n-m)x$,在整周期上的积分是零.

热棒的区间是 $[0,\ell]$.令 $\theta=\pi x/\ell$,则 $\mathrm{d}x=(\ell/\pi)\,\mathrm{d}\theta$,$x$ 从 $0$ 到 $\ell$ 时 $\theta$ 从 $0$ 到 $\pi$.
\begin{equation}
  \int_0^{\ell}
    \sin\Bigl(\frac{n\pi x}{\ell}\Bigr)
    \sin\Bigl(\frac{m\pi x}{\ell}\Bigr)
    \,\mathrm{d}x
  =\frac{\ell}{\pi}\int_0^{\pi}\sin(n\theta)\sin(m\theta)\,\mathrm{d}\theta
  =\frac{\ell}{2}\,\delta_{mn}.
\end{equation}
$[0,\pi]$ 上的积分是 $[0,2\pi]$ 上的一半,所以范数从 $\pi$ 变成 $\ell/2$.用 $\sin(m\pi x/\ell)$ 去乘正弦级数再积分,只留下 $A_m\cdot\ell/2$.
\end{solution}

因此振幅系数
\begin{equation}
  A_n = \frac{2}{\ell} \int_0^{\ell} f(x) \sin{ \left( \frac{n\pi}{\ell} x \right) } dx 
\end{equation}
$A_n$称为傅里叶级数的\textbf{系数}.若所有系数都求得后,我们就得到了$f(x)$.原则上,所有的平滑的单值函数并没有什么问题,
但似乎缺乏数学的严谨.自然,越多的级数项对函数的逼近也就越好.

\subsection{傅里叶定理}
\label{subsec:fourier_theorem}
对于任意一个以$2\ell$为周期的函数,
$$
   f(x + 2\ell) = f(x),  
$$
我们可以通过三角函数族进行级数展开
\begin{equation}
  f(x) = \frac{a_0}{2} + \sum_{n=1}^{\infty} \left[ a_n \cos{ \frac{n\pi}{\ell} x } + b_n \sin{ \frac{n\pi}{\ell} x } \right] 
\end{equation}
利用三角函数族的正交性,可以得到展开系数
\begin{align}
  a_n = \frac{1}{\ell} \int_{-\ell}^{\ell} f(x) \cos {  \left( \frac{n\pi}{\ell} x \right) } dx
  \\
  b_n = \frac{1}{\ell} \int_{-\ell}^{\ell} f(x) \sin {  \left( \frac{n\pi}{\ell} x \right) } dx
\end{align}
这里的积分上下限为一个周期$2\ell$,也可以为$[0,2\ell]$.常数项写成 $a_0/2$,是为了让 $n=0$ 仍用同一个积分.$\int_{-\ell}^{\ell}1\,\mathrm{d}x=2\ell$,而 $n\geq 1$ 时 $\int_{-\ell}^{\ell}\cos^2(n\pi x/\ell)\,\mathrm{d}x=\ell$.展开式乘 $1$ 再积分,左边是 $\int_{-\ell}^{\ell}f\,\mathrm{d}x$,右边是 $(a_0/2)\cdot 2\ell$,所以 $a_0$ 正好等于 $n=0$ 的余弦公式.级数里的常数项 $a_0/2$ 就是 $f$ 在一个周期上的平均值.
上述傅里叶级数展开的成立条件被称为\textbf{狄里希利条件} (Dirichlet conditions),即$f(x)$在$\left[-\ell, \ell\right]$区间内只有有限个间断点,且每个周期内有有限个极值点.满足
这两个条件的函数称为\textbf{分段平常}.
满足狄里希利条件的函数$f(x)$在点$x_0$处间断,那么其傅里叶级数在此点的值为该函数左右值的算术平均
\begin{equation}
  \textrm{级数和}(x_0) = \lim_{\epsilon \to 0} \half \left[
     f(x_0 + \epsilon) + f(x_0  - \epsilon) \right]
\end{equation}

对于奇函数和偶函数的傅里叶展开,不难发现,奇函数的展开为

\begin{equation}
  f(x) = \sum_{n=1}^{\infty} b_n \sin {  \left( \frac{n\pi}{\ell} x \right) }
\end{equation}
称为\textbf{傅里叶正弦级数}.
而偶函数的展开为
\begin{equation}
  f(x) = \frac{a_0}{2} + \sum_{n=1}^{\infty}  a_n \cos{  \left( \frac{n\pi}{\ell} x \right) }
\end{equation}
称为\textbf{傅里叶余弦级数}.

只在 $(0,\ell)$ 上给定的函数可以有很多种延拓,它们在 $(0,\ell)$ 上都代表原来的 $f$.边界条件常常把延拓定下来.要求 $f(0)=f(\ell)=0$ 时,把 $f$ 延拓成以 $2\ell$ 为周期的奇函数:在 $(0,\ell)$ 上 $g=f$,而 $-\ell<x<0$ 时 $g(x)=-f(-x)$.由周期性,$g(\ell+\epsilon)=g(-\ell+\epsilon)$;再由奇性,$g(-\ell+\epsilon)=-g(\ell-\epsilon)$.因此 $g$ 关于 $x=\ell$ 也是奇的.连续延拓只能有 $g(0)=g(\ell)=0$.正弦级数在这两个端点自动为零,系数是
\begin{equation}
  b_n = \frac{2}{\ell}\int_0^{\ell} f(x)\sin\Bigl(\frac{n\pi x}{\ell}\Bigr)\,\mathrm{d}x.
\end{equation}
热棒两端保持零度,用的就是这个奇延拓.

要求 $f'(0)=f'(\ell)=0$ 时,改成偶延拓,$g(-x)=g(x)$.偶函数的导数是奇函数,左右导数在 $x=0$ 处相等当且仅当 $f'(0)=0$.把原点移到 $x=\ell$,同样得到 $f'(\ell)=0$.余弦级数的导数是正弦级数,在两端自动为零.系数是
\begin{equation}
  a_n = \frac{2}{\ell}\int_0^{\ell} f(x)\cos\Bigl(\frac{n\pi x}{\ell}\Bigr)\,\mathrm{d}x.
\end{equation}
只在 $(0,\ell)$ 上给定、再选一种延拓,得到的正弦级数或余弦级数称为\idx{半幅级数}{banfujishu}.同一定义可以有两个级数,它们在 $(0,\ell)$ 内都等于 $f$,在区间外代表不同的延拓.

\begin{example}
$f(x)=x$ 定义在 $(0,\ell)$ 上.分别求它的正弦级数和余弦级数.
\end{example}
\begin{solution}
奇延拓后 $b_n=(2/\ell)\int_0^{\ell}x\sin(n\pi x/\ell)\,\mathrm{d}x$.分部积分,$\int_0^{\ell}\cos(n\pi x/\ell)\,\mathrm{d}x=0$,只剩边界项,
\begin{equation}
  b_n=\frac{2\ell}{n\pi}(-1)^{n+1},
  \qquad
  x=\frac{2\ell}{\pi}\sum_{n=1}^{\infty}\frac{(-1)^{n+1}}{n}\sin\Bigl(\frac{n\pi x}{\ell}\Bigr),
  \quad 0<x<\ell.
\end{equation}
端点 $x=0$ 和 $x=\ell$ 上级数和是 $0$,与 $f$ 在端点的值一致.把这个奇延拓写到 $(-\ell,\ell)$ 上,就是后面的锯齿函数.

偶延拓是 $|x|$.$a_0=(2/\ell)\int_0^{\ell}x\,\mathrm{d}x=\ell$,常数项是 $\ell/2$.$n\geq 1$ 时分部积分,
\begin{equation}
  \int_0^{\ell}x\cos\Bigl(\frac{n\pi x}{\ell}\Bigr)\,\mathrm{d}x
  =\frac{\ell^2}{n^2\pi^2}\bigl[(-1)^n-1\bigr],
\end{equation}
边界上的正弦项是零.因此偶数 $n$ 的 $a_n$ 是零,奇数 $n=2k+1$ 给出 $a_n=-4\ell/(n^2\pi^2)$,
\begin{equation}
  x=\frac{\ell}{2}-\frac{4\ell}{\pi^2}\sum_{k=0}^{\infty}
    \frac{1}{(2k+1)^2}\cos\Bigl(\frac{(2k+1)\pi x}{\ell}\Bigr),
  \quad 0<x<\ell.
\end{equation}
令 $x=0$,左边是 $0$,得到 $\sum_{k=0}^{\infty}1/(2k+1)^2=\pi^2/8$.
\end{solution}

\begin{example}
$f(x)=e^{ax}$ ($a\neq 0$) 在 $(-\pi,\pi)$ 上的全级数.周期取 $2\pi$,即 $\ell=\pi$.
\end{example}
\begin{solution}
\begin{equation}
  a_0=\frac{1}{\pi}\int_{-\pi}^{\pi}e^{ax}\,\mathrm{d}x=\frac{2\sinh(a\pi)}{a\pi},
\end{equation}
常数项是 $a_0/2=\sinh(a\pi)/(a\pi)$.对 $n\geq 1$,
\begin{equation}
  \int e^{ax}\cos(nx)\,\mathrm{d}x
  =\frac{e^{ax}\bigl(a\cos(nx)+n\sin(nx)\bigr)}{a^2+n^2}.
\end{equation}
$\sin(n\pi)=0$,$\cos(n\pi)=(-1)^n$,从 $-\pi$ 积到 $\pi$ 得到 $2a(-1)^n\sinh(a\pi)/(a^2+n^2)$.除以 $\pi$,
\begin{equation}
  a_n=\frac{2a(-1)^n\sinh(a\pi)}{\pi(a^2+n^2)}.
\end{equation}
正弦积分
\begin{equation}
  \int e^{ax}\sin(nx)\,\mathrm{d}x
  =\frac{e^{ax}\bigl(a\sin(nx)-n\cos(nx)\bigr)}{a^2+n^2}
\end{equation}
同样只留下余弦边界项,
\begin{equation}
  b_n=\frac{2n(-1)^{n+1}\sinh(a\pi)}{\pi(a^2+n^2)}.
\end{equation}
$a_n$ 和 $b_n$ 都不为零:指数函数既不是奇函数也不是偶函数.系数随 $1/n^2$ 和 $1/n$ 衰减,间断在周期端点 $x=\pm\pi$,那里级数和是两端极限的平均 $\cosh(a\pi)$.
\end{solution}

对于复指数的展开,不难发现可以写成
\begin{equation}
  f(x) = \sum_{n=-\infty}^{\infty} c_n e^{\imath \frac{n\pi}{\ell} x},
\end{equation}
其中
$c_n = \half(a_n - \imath b_n)$, $c_{-n} = \half (a_n + \imath b_n)$, $n>0$, $c_0 = \half a_0$.
写出来为
\begin{equation}
  c_n = \frac{1}{2\ell} \int_{-\ell}^{\ell} f(x) e^{-\imath \frac{n\pi}{\ell} x} dx ,
\end{equation}
尽管$f(x)$是实数,但其傅里叶系数却可能是复数,还可以看出$c_{-n} = c_{n}^*$.
复指数函数族也是正交的
$$
  \braket{ e^{\imath \frac{n\pi}{\ell} x } | e^{\imath \frac{m\pi}{\ell} x } } = 2\ell \delta_{mn} .
$$

\begin{example}
对于锯齿函数
$$  
f(x)= \begin{cases}
  x,   &  0 < x \leq  \ell 
  \\
  x - 2 \ell,   & \ell < x \leq 2\ell
\end{cases}
$$
求其傅里叶级数.
\end{example}
\begin{solution}
把一个周期挪到 $(-\ell,\ell)$,这个函数就是奇函数 $f(x)=x$,间断点在 $\pm\ell$.因此 $a_n=0$,并且
\begin{equation}
  b_n = \frac{2}{\ell}\int_0^{\ell} x\sin\Bigl(\frac{n\pi x}{\ell}\Bigr)\,\mathrm{d}x .
\end{equation}
分部积分,令 $u=x$,$\mathrm{d}v=\sin(n\pi x/\ell)\,\mathrm{d}x$,边界项是 $-\ell^2(-1)^n/(n\pi)$,剩余的余弦积分等于零,所以
\begin{equation}
  b_n = \frac{2\ell}{n\pi}(-1)^{n+1}.
\end{equation}
因此,该级数为
\begin{align}
f(x) &=  \sum_{n=1}^{\infty} \frac{2 \ell }{n\pi} (-1)^{n+1} \sin {  \left( \frac{n\pi}{\ell} x \right) }  \nonumber
\\
 & =\frac{2\ell}{\pi} \left[ \sin{\frac{\pi x}{\ell} } - \frac{1}{2}\sin{\frac{2\pi x}{\ell} } 
  +  \frac{1}{3}\sin{\frac{3\pi x}{\ell} } + \cdots + \frac{(-1)^{n+1}}{n}\sin{\frac{n\pi x}{\ell} }
 \right] \nonumber
.
\end{align}
容易验证
$$
  f(0) = 0; f(\ell) = 0;
$$
当$x=\ell/2$, 
$$
f(\ell/2) = \ell/2 =   \frac{2\ell}{\pi} \left[ 1 - 0 - \frac{1}{3} + \frac{1}{5} - 0 -\frac{1}{7}+ \cdots \right]
$$
因此,我们得到莱布尼兹等式,即以下等式
$$
\frac{\pi}{4} = 1-\frac{1}{3} + \frac{1}{5} - \frac{1}{7} + \cdots = \sum_{n=0} \frac{(-1)^n}{2n + 1} .
$$
上式可以通过等式
$\int_0^1 \frac{d x}{1+x^2}=\left.\tan ^{-1} x\right|_0 ^1=\frac{\pi}{4}$,和被积函数的级数展开验证.
其实,这里的锯齿函数可以直接用$f(x) = x, -\ell < x < \ell$来表示,相应的间断点则被移到$\pm \ell$处.
\end{solution}


\subsection{能量定理}
\idx{帕塞瓦尔等式}{paseiwaer}
由函数的复数傅里叶级数展开可以得到
\begin{equation}
  E = \frac{1}{2\ell} \int_{-\ell}^{\ell} dx f(x)^* f(x) 
 = \frac{1}{2\ell} \int_{-\ell}^{\ell} dx  \sum c_n e^{\imath \frac{n\pi}{\ell} x} \sum c_m e^{-\imath \frac{m\pi}{\ell} x}
 = \sum |c_n|^2
\end{equation}
利用前面的等式$c_n = \half(a_n - \imath b_n)$,  $c_{-n} = \half (a_n + \imath b_n)$, $n>0$, $c_0 = \half a_0$ 可得
\begin{equation}
  E = \frac{a_0^2}{4}+\sum_{n=1}^{\infty}\left(\frac{a_n^2}{2}+\frac{b_n^2}{2}\right) .
\end{equation}
称为\textbf{Parseval等式}(Parseval's identity).
其物理意义指的是,一复杂波形分解为许多波幅为$c_n$的正(余)弦波的叠加,该波所携带的总能量为各个成分波强度的求和.

\begin{example}
矩形波.取 $\ell=1$,
\begin{equation}
  f(x)=\begin{cases}
    -1, & -1<x<0,\\
    +1, & 0<x<1,
  \end{cases}
\end{equation}
再以 $2$ 为周期延拓.
\end{example}
\begin{solution}
$f$ 是奇函数,$a_n=0$.
\begin{equation}
  b_n=2\int_0^1\sin(n\pi x)\,\mathrm{d}x
  =\frac{2\bigl(1-(-1)^n\bigr)}{n\pi}.
\end{equation}
偶数 $n$ 的系数是零,奇数 $n=2k+1$ 的系数是 $4/(n\pi)$,
\begin{equation}
  f(x)=\frac{4}{\pi}\sum_{k=0}^{\infty}
    \frac{\sin\bigl((2k+1)\pi x\bigr)}{2k+1}.
\end{equation}
$x=0$ 是跳跃点,左右极限是 $-1$ 和 $1$,级数和是它们的平均 $0$,每一项在这里确实是零.$x=1/2$ 处 $f=1$,代入得到莱布尼兹级数 $\pi/4=1-1/3+1/5-\cdots$.跳跃附近部分和会冲过阶跃一段固定的比例,再在两侧振荡.项数增加时,这段过冲并不变小,只是挤向间断点.平均平方误差仍然趋于零:逐点图像和平均平方逼近回答的是两件不同的事.

$|f|^2=1$,所以 $E=1/(2\ell)\int_{-\ell}^{\ell}|f|^2\,\mathrm{d}x=1$.Parseval 等式只留下奇数项,
\begin{equation}
  1=\sum_{k=0}^{\infty}\frac12\Bigl(\frac{4}{(2k+1)\pi}\Bigr)^2
  =\frac{8}{\pi^2}\sum_{k=0}^{\infty}\frac{1}{(2k+1)^2}.
\end{equation}
因此 $\sum_{k=0}^{\infty}1/(2k+1)^2=\pi^2/8$.这和半幅余弦级数在 $x=0$ 处得到的是同一个和.
\end{solution}

\begin{example}
$f(x)=x^2$ 在 $[-\pi,\pi]$ 上,周期 $2\pi$.
\end{example}
\begin{solution}
这是偶函数,$b_n=0$.
\begin{equation}
  a_0=\frac{1}{\pi}\int_{-\pi}^{\pi}x^2\,\mathrm{d}x=\frac{2\pi^2}{3},
\end{equation}
常数项 $a_0/2=\pi^2/3$.分部积分两次,边界上的 $\sin(nx)$ 是零,
\begin{equation}
  \int_0^{\pi}x^2\cos(nx)\,\mathrm{d}x=\frac{2\pi(-1)^n}{n^2},
  \qquad
  a_n=\frac{2}{\pi}\cdot\frac{2\pi(-1)^n}{n^2}=\frac{4(-1)^n}{n^2}.
\end{equation}
\begin{equation}
  x^2=\frac{\pi^2}{3}+4\sum_{n=1}^{\infty}\frac{(-1)^n}{n^2}\cos(nx),
  \qquad -\pi\leq x\leq\pi.
\end{equation}
$x=\pi$ 处函数连续,$x^2=\pi^2$,$\cos(n\pi)=(-1)^n$,
\begin{equation}
  \pi^2=\frac{\pi^2}{3}+4\sum_{n=1}^{\infty}\frac{1}{n^2},
\end{equation}
所以 $\sum_{n=1}^{\infty}1/n^2=\pi^2/6$.

再用 Parseval 等式核对.左边
\begin{equation}
  E=\frac{1}{2\pi}\int_{-\pi}^{\pi}x^4\,\mathrm{d}x=\frac{\pi^4}{5}.
\end{equation}
右边是 $a_0^2/4+\sum a_n^2/2=(2\pi^2/3)^2/4+\sum_{n=1}^{\infty}\frac12\bigl(4/n^2\bigr)^2=\pi^4/9+8\sum 1/n^4$.这个等式给出的是 $\sum 1/n^4$,不是 $\sum 1/n^2$.$\pi^2/6$ 来自把级数代回 $x=\pi$,两边各自算完再比较.
\end{solution}

\subsection{三角系的完备性}
前面把完备说成:平均平方误差可以小于任意正数.这里对分段平常的周期函数把这件事证完.先作换元 $\theta=\pi x/\ell$,把区间 $[-\ell,\ell]$ 化到 $[-\pi,\pi]$.两边的积分只差常数 $\ell/\pi$,所以只要在 $[-\pi,\pi]$ 上证明即可.此时
\begin{equation}
  a_n = \frac{1}{\pi}\int_{-\pi}^{\pi} f(x)\cos(nx)\,\mathrm{d}x,
  \qquad
  b_n = \frac{1}{\pi}\int_{-\pi}^{\pi} f(x)\sin(nx)\,\mathrm{d}x.
\end{equation}

设 $s_n$ 是 $f$ 的 $n$ 阶傅里叶部分和,$T$ 是次数不超过 $n$ 的三角多项式,$r=T-s_n$.由系数定义,$\int_{-\pi}^{\pi}(f-s_n)\cos(kx)\,\mathrm{d}x=0$,正弦同样,所以 $\int(f-s_n)r\,\mathrm{d}x=0$.于是
\begin{equation}
  \int_{-\pi}^{\pi}(f-T)^2\,\mathrm{d}x
  = \int_{-\pi}^{\pi}(f-s_n)^2\,\mathrm{d}x + \int_{-\pi}^{\pi} r^2\,\mathrm{d}x.
\end{equation}
第二项非负,因此在这个次数之内,$s_n$ 的平均平方误差最小.把平方展开,并用 $\int_{-\pi}^{\pi}\cos^2(kx)\,\mathrm{d}x=\pi$($k\geq 1$)和 $\int_{-\pi}^{\pi}1\,\mathrm{d}x=2\pi$,得到
\begin{equation}
  \int_{-\pi}^{\pi}(f-s_n)^2\,\mathrm{d}x
  = \int_{-\pi}^{\pi} f^2\,\mathrm{d}x
    - \pi\Biggl(\frac{a_0^2}{2}+\sum_{k=1}^{n}(a_k^2+b_k^2)\Biggr).
\end{equation}
左边非负,这就是三角系上的贝塞尔不等式.要证明完备,还要造出一组三角多项式,把这个差压到任意小.

部分和可以写成卷积.记\textbf{狄里希利核}
\begin{equation}
  D_n(\theta) = \half + \sum_{k=1}^{n}\cos(k\theta)
  = \frac{\sin\bigl((n+\half)\theta\bigr)}{2\sin(\theta/2)}.
\end{equation}
第二个等号用 $2\sin(\theta/2)$ 去乘求和式,积化和差后中间项相消.于是
\begin{equation}
  s_n(x) = \frac{1}{\pi}\int_{-\pi}^{\pi} f(t)\,D_n(t-x)\,\mathrm{d}t,
\end{equation}
并且 $\int_{-\pi}^{\pi} D_n(\theta)\,\mathrm{d}\theta=\pi$.

把前 $n+1$ 个狄里希利核取平均,得到\textbf{费耶核}
\begin{equation}
  F_n(\theta)
  = \frac{1}{n+1}\sum_{m=0}^{n} D_m(\theta)
  = \frac{1}{2(n+1)}
    \Biggl(\frac{\sin((n+1)\theta/2)}{\sin(\theta/2)}\Biggr)^2.
\end{equation}
它非负,积分仍是 $\pi$.对应的\textbf{费耶平均} $\sigma_n=(s_0+\cdots+s_n)/(n+1)$ 仍是次数不超过 $n$ 的三角多项式,并且
\begin{equation}
  \sigma_n(x)-f(x)
  = \frac{1}{\pi}\int_{-\pi}^{\pi}\bigl[f(x-t)-f(x)\bigr] F_n(t)\,\mathrm{d}t.
\end{equation}

先设 $f$ 连续,并且 $f(-\pi)=f(\pi)$,从而在圆周上一致连续.给定 $\epsilon>0$,取 $\delta$ 使 $|t|<\delta$ 时 $|f(x-t)-f(x)|<\epsilon$.积分拆成 $|t|<\delta$ 和 $|t|\geq\delta$.$F_n\geq 0$,峰内的贡献小于 $\epsilon$.峰外 $|\sin(t/2)|\geq\sin(\delta/2)$,所以 $F_n$ 不超过 $1/\bigl(2(n+1)\sin^2(\delta/2)\bigr)$,$n$ 大时峰外的贡献可以小于 $\epsilon$.因此 $\sigma_n$ 一致收敛到 $f$,
\begin{equation}
  \int_{-\pi}^{\pi}(f-\sigma_n)^2\,\mathrm{d}x
\end{equation}
可以任意小.$s_n$ 的平均平方误差不超过 $\sigma_n$ 的,所以连续且周期闭合的函数已经可以用三角多项式在均方意义下逼近.

若 $f$ 只是分段平常,一个周期里只有有限个跳跃.在每个跳跃处用一段宽度为 $\eta$ 的直线把左右值接上,端点 $f(-\pi)$ 与 $f(\pi)$ 同样接上,得到连续的周期函数 $h$.这些小区间上 $|f-h|$ 有界,$\eta$ 足够小时 $\int(f-h)^2$ 可以任意小.再对 $h$ 取三角多项式 $T$,使 $\int(h-T)^2$ 任意小.由 $(A+B)^2\leq 2A^2+2B^2$,
\begin{equation}
  \int(f-T)^2\,\mathrm{d}x
  \leq 2\int(f-h)^2\,\mathrm{d}x + 2\int(h-T)^2\,\mathrm{d}x
\end{equation}
也可以任意小.$n$ 不小于 $T$ 的次数时,$\int(f-s_n)^2\leq\int(f-T)^2$.平均平方误差可以任意小,三角函数族完备.回到 $[-\ell,\ell]$,贝塞尔不等式里的余项趋于零,就得到正文中的 Parseval 等式.狄里希利条件另外保证间断点处级数和等于左右值的平均;这里证明的是平均平方意义下的完备.